
\documentclass[]{article}

\usepackage{booktabs}
\usepackage[russian]{babel}

\begin{document}




I regard these variables as possible predictors rather than factors of reporting \citep{wikstrom2011does}...

\cite{goudriaan2004reporting} juxtapose rational cost-benefit calculus and "normative response": \textit{ In contrast, normative responses are not directly caused by expected costs or benefits, but by norms that exist in the victim's social context. "Crimes should be reported to the police," "this is not a case for the police," or "I should deal with this myself' are examples of possible norms that may influence crime reporting.}. This seems a highly doubtful analysis.

A visual summary of causal relations in my model provided in Fig \ref{fig:causal_diagram} makes it evident that an integrative framework is not merely reasonable or theoretically sound but plain necessary if one aims to identify the effect of social support on reporting. Almost all the factors that are supposed to influence the decision to notify the police may as well influence the decision to seek help in friends and family. In other words, the dependent and the focus independent variables have multiple confounders. This poses a threat to identification, which I am able to overcome by inclusion of a broad range of controls. However, there are three caveats. 

First, I do not observe a victim's cultural context. The only way I am able to try to control for it is to introduce federal district fixed effects and settlement type variable, which is better than nothing but admittedly a very crude solution,

Second, I do not observe offenders characteristics other than their gender and relation to a victim. I do not include offender gender or interaction of victim-offender genders in my modeling because in my data almost the crime was committed by men. I do include relationship distance which proves to be an important factor. Still, I know nothing of offender's resourcefulness which must be one of the serious considerations in a victim's decision making process. This is an important omission, but one that my study shares with virtually all the research on reporting.
\textit{TBAa: fear of reprisal \citep{felson1999victim}}

Third, I know only if the police was notified and if a victim told about the incident to someone in their intimate circle -- and I do not know which came first. This means endogeneity is possible: a victim may report to the police and then relate the whole story to friends or family \textit{because} of how the interaction with the police has worked out. While our knowledge of post-victimization behavior suggests that if someone is notified at all then friends or family come before the police \citep{ruback2001social, knoth2019reporting}, the doubt remains. I address this issue in Subsection \ref{subsec:robust} by introducing a supporting model of different design. 

Motivation for year fixed effects:
\textit{Some scholars have attributed substantive meanings to this change and argued that it was caused by changes in sex-role attitudes in the United States (e.g., Orcutt \& Faison 1988) or by the success of rape law reforms and media campaigns (e.g., Bachman 1993, Spohn \& Horney 1992). The evidence from the United Kingdom similarly supported the notion that research and policy attention to rape, sexual assault, and domestic violence had led to more adequate handling of such crimes and that this had resulted in a greater willingness of victims to report crimes by intimate friends and relatives (Tarling \& Morris 2010).}





Traditionally \citep{skogan1976citizen, skogan1984reporting, gottfredson1979study, tarling2010reporting}, gravity of an incident, i.e. perceived actual and potential harm, is seen as the main predictor of reporting to the police. Relational distance between victim and offender, especially in violent crime, has also received much attention from researchers and was found to be an important factor in post-victimization behavior \citep{felson1999victim, kaukinen2002help, knoth2019reporting}. 



Recent study by \cite{dzmitryieva2025crime} proved these predictors to be relevant in Russian context, too.


\footnote{victimization survey data that I use features only one question that captures social interaction that follows victimization and, arguably, precedes notifying the police: respondents who said they had experienced some kind of victimization were asked if they had told anyone in their intimate circle about the incident.}

 I focus on turning to friends or family -- the event that, I speculate, must be common to most of post-victimization help seeking paths. An individual either activates their network or they cope on their own.
 
 	
 \textit{TBA: add on \cite{pescosolido1998social}.} In this study I am making a step in this direction. There are several paths one can take in seeking help, some of which lead to or pass through notifying the police. For example, a person can inform a friend who suggests to visit a crisis center where, among other things, the person is recommended to file a case against the offender. It is these paths and the factors of them going through or passing round the police station that need to be studied -- especially if we are interested in applicable results. Unable to observe these paths in full due to data limitations (see the Section \ref{sec:data}), in this study I focus on turning to friends or family -- the event that, I speculate, must be common to most of post-victimization help seeking paths. An individual either activates their network or they cope on their own. In the former case the shock spread triggers an essentially collective response to the cause: not only the victim receives advice, emotional and practical support (or, sadly, further trauma), not only other parties become involved into decision making process -- sometimes they act on behalf of the victim, either notifying the police, seeking informal help, or resorting to self-help tactics.

		Recency problem \citep{goudriaan2004reporting}):
\textit{The use of a 5-year reference period has probably resulted in an underestimate of victimization rates and an overestimate of the proportion reported to the police. Memory decays with time and the longer the reference period, the less complete the recall and reporting of crime events. Because events reported to the police are rehearsed and have consequences subsequent to the event, they will be remembered more frequently than events that are not reported to the police. Biderman and Lynch (1981) have shown, however, that this type of recency bias is not correlated with other attributes of the crime, so it is unlikely to affect the analyses presented here. They compared attributes of events reported in the most recent month of the reference period for the National Crime Survey with those reported in the most distant month of the reference period. The distributions of these event attributes were similar for the most proximate and the most distant month of the reference period for all attributes except reporting to the police. The percent reported to the police in the most distant month was significantly greater than the proportion reported to the police in the most proximate month.}

\textit{data limitations}
The RCVS includes almost no questions regarding sex-related crimes and intimate partner violence (IPV), which likely leads to a significant underestimation of the overall victimization rate in the country.

In the late 2010s, cybercrimes became the most frequently reported offenses by respondents in Russia 30. We hypothesize that victims of cybercrime are more prone to providing false positive (Type I error) or false negative (Type II error) responses regarding their victimization experience. In other words, due to the changing structure of crime, more respondents might identify as victims of cybercrime when, in reality, they were targets of attempted crimes without incurring property damage. Conversely, more respondents might decline to report actual victimization, considering it not ‘noteworthy’ enough to mention.

updated regulations came into force only after the fieldwork had been completed.  A further methodological risk lies in the non-obvious correlation between the probability of initial contact or probability of cooperation and the likelihood of having experienced cybercrime victimization (much of which occurs via telephone or online messengers). This suggests a potential self-selection bias, which warrants further investigation.

According to CMLC data, 3.7\% (95\% CI: 3.6–3.8\%) of the Russian population were a victim of a crime in 2024. The comparable RCVS estimate is 6.4\%. However, after excluding cybercrimes and cybercrime attempts, the RCVS victimization rate drops to 4.1\% (95\% CI: 3.7–4.4\%). This discrepancy arises because the CMLC’s closed-ended list of victimization experiences does not include cybercrime.



Thus, two methodologically distinct surveys, with different samples and modes of data collection (computer-assisted telephone Interviewing vs. face-to-face), yield equivalent victimization estimates, and sample estimates were verified both with and without weighting.



A by-product of my work is an analysis of personal covariates of seeking help in intimate circle ...

The rest of the paper is conveniently organized the way you would expect. After the literature review in Section \ref{sec:lit_rev}, I will introduce my data (Section \ref{sec:data}) and analytical approach (Section \ref{sec:approach}), present my results (Section \ref{sec:results}), tackle possible effect identification problems (Subsection \ref{subsec:robust}), offer a discussion of findings and the methodological changes they call for (Section \ref{sec:discussion}), and conclude with conclusion (Section \ref{sec:conclusion}).









Ecological model of reporting did not dethrone the cost-benefit view of victim's decision making. In part, because 


Rather, it provided correction by taking into account environment in which calculus takes place. What \cite{pescosolido1992beyond} wrote about the revised rational choice perspective on how people seek medical help is also right about the ecological approach in reporting literature: it still `focuses on the individual but tilts the analytic lens upward in an attempt to capture other larger conditions (including networks) that affect the individual unit of analysis'. 


Despite \cite{goudriaan2004reporting} did not find significant connection between reporting and several meso- and macro-level characteristics available in ICVS, other researchers followed suite (which tells something about perceived soundness of the ecological metaphor).  







\textit{TBA: examples of studies employing ecological framework and their findings \citep{berg2017mobilization, baumer2002neighborhood, goudriaan2006neighbourhood}.}


	
\textit{Mixed results:	For example, using the National Crime Victimization Survey, Baumer (2002) found that the reporting of simple assault generally increased with neighborhood poverty except at high levels of poverty, where the probability of reporting dropped sharply. In comparison, Goudriaan, Wittebrood, and Nieuwbeerta (2006) observed a negative relationship between crime reporting and neighborhood poverty, but like Baumer they also observed a sharp drop in crime reporting at high levels of poverty.}

\textit{or example, Goudriaan, Wittebrood, and Nieuwbeerta (2006) found that neighborhood mobility was related negatively to victimization reporting and that the effect of this variable was mediated by neighborhood social cohesion and confidence in the police.}


	\textit{TBA: on attitudes literature \citep{boateng2018crime,ann2010neighborhood,warner2007directly} and victim-offender relationship literature: }



\cite{bruch2017decision} :

\textit{Quantitative data analysis, despite focusing primarily on individual-level outcomes, is typically conducted without any reference to—let alone a model of—individual action (Goldthorpe 1996; Esser 1996). Rational choice provides a theory of action that can anchor empirical research in meaningful descriptions of individuals’ behavior (Hedström \& Swedberg 1996). Importantly, the choice behavior of rational actors can also be straightforwardly implemented in regression-based models readily available in statistical software packages.}

\textit{Most relevant to this article, sociologists have argued that the focus in classical rational choice on the individual as the primary unit of decision-making represents a fundamentally asocial representation of behavior. In moving beyond rational choice, theories of decision-making in sociology highlight the importance of social interactions and relationships in shaping behavior (Pescosolido 1992; Emirbayer 1997).}

\textit{People prefer not to make these emotionally difficult tradeoffs, and that shapes decision strategy selection (Hogarth 1991; Baron 1986; Baron \& Spranca 1997). Experiments on these types of emotional decisions have shown that, when facing emotionally difficult decisions, decision-makers avoid compensatory evaluation and instead select the alternative that is most attractive on whatever dimension is difficult to trade off (Luce et al.2001; Luce et al.1999). Thus, the emotional valence of specific options shapes decision strategies.}

\textit{This has led to a large body of work on the affect heuristic, which is grounded in the idea that people have positive and negative associations with different stimuli, and they consult this “affect pool” when making judgments. This shortcut is often more efficient and easier than cognitive strategies such as weighing pros and cons or even disjunctive rules for evaluating the relative merits of each choice outcome (Slovic etal.2004). Affectparticularly how it relates to decision-making—is rooted in dual process accounts of human behavior. The basic idea is that people experience the world in two different ways: one that is fast, intuitive, automatic, and unconscious, and another that is slow, analytical, deliberate, and verbal (Evans 2008; Kahneman 2011). A defining characteristic of the intuitive, automatic system is its affective basis (Epstein 1994). Indeed, affective reactions to stimuli are often the very first reactions people have. Having determined what is salient in a given situation, affect thus guides subsequent processes, such as information processing, that are central to cognition (Zajonc 1980).}

\textit{Choice outcomes are also highly influenced by what option is identified as the “default.” Defaults are whatever happens in a decision if the chooser “decides not to decide” (Feinberg \& Huber 1996). Defaults exert a strong effect on people’s choices, even when the stakes of the decision are high (Johnson etal.1993; Johnson \& Goldstein 2013). A number of sociological studies have theorized about how people “drift” into particular outcomes or situations (e.g., Matza 1967); defaults exist in part because choices are embedded in specific environments that emphasize one set of options over others.8}





***

\cite{thoits1995stress}:
\textit{Major Findings and Gaps: Coping Strategies. Turning to research on copingstrategiesa, numberof studiesindicatethatindividualstypicallyuse multipletacticswhencopingwith majorlifeeventsor ongoingstrains(e.g., Billings,Cronkite,and Moos 1983; Folkmanand Lazarus 1980;StoneandNeale 1984). FolkmanandLazarus(1980) reportedthatin98 percent of 1,300 stressfuelpisodes,subjectsused bothproblem-focuseadnd emotion-focusecdoping strategiesT.hesefindingssupportetdheircontentiotnhatthereareusuallytwosourcesofstress whichmustbe handled,bothsituationaldemandsand one's emotionalresponseto those demands. Not surprisingl syt,ressorswhichare appraisedas moresevere(e.g., as harms/losseosr as threats)evokegreaternumbersof copingresponses(Cronkiteand Moos 1984; Folkmanand Lazarus 1980;McCrae 1984;Menaghan1982, 1983;MenaghanandMerves1984). In general, problem-focusecdopingis morelikelywhensituationadlemandsareappraisedas controllable; emotion-focusecdopingis morelikelywhendemandsseem uncontrollabl(eBillingset al. 1983; Coyne,Aldwin,and Lazarus 1981; Folkman1984; Folkmanand Lazarus 1980, 1985; Folkmanet al. 1986; Forsytheand Compas 1987; Stone and Neale 1984; Thoits 1991). Althoughthe perceiveduncontrollabilitoyf a stressorconsistentlypredictsthe use of emotion-focusesdtrategiesi,t shouldbe notedthatsome studiesfindno associationbetween perceivedsituationaclontrolandefforttsoproblem-solv(eStoneandNeale 1984;Thoits1991). Coping researchersgenerallyexpectproblem-focusecdoping to be more beneficialfor}

\textit{Men moreoftenreportcontrollintgheiremotions, acceptingthe problem,not thinkingaboutthe situation,and engagingin problem-solving efforts.Womenmoreoftenreportseekingsocial support,distractintghemselvesl,ettingout theirfeelings,and turningto prayer.Women's greaterpropensityto seek social supportis especiallyconsistenatcrossstudies.Buttherearea numberofexceptionsintheliteraturwe ith respectto genderdifferenceisn problem-focusecdoping (e.g., Billings and Moos 1984; FolkmanandLazarus 1980; Ross andMirowsky1989; MilkieandThoits1993). This maybe because men's and women's use of problem-focusecdopingmay dependupon perceiving controlorpowerin a roledomain-forexample,menintheoccupationalarenaandwomenin thefamilyarena(Folkmanand Lazarus 1980; Menaghan1982; Pearlinand Schooler1978).}

\textit{.The effectsof perceived social supporthave mostfrequentlbyeen examinedin theliteraturee,speciallytheeffectsof perceivedemotionalsupport(i.e., beliefsthatlove and caring,sympathyand understanding, and/oresteemand value are available fromsignificanotthers).The perceptionor beliefthat emotionalsupportis available appearsto be a muchstrongeirnfluenceon mentalhealththan the actual receiptof social support(Dunkel-Schetterand Bennett1990; Wethingtonand Kessler 1986).}

***






Scarcity of research on connection between seeking informal help and reporting may in part be due to there are seemingly no immediate policy implications \citep{greenberg1992after}. However, 1) it may be more effective to work through situational advisors (victim decision making being compromised by shock), i.e. potentially anyone; 2) if we believe that consulting others heightens probability of reporting we must conclude that lonely people report less, which gives us an insight on the direction of bias introduced by under-reporting crime.

\cite{pescosolido1992beyond} 1) victimization is a shock to the whole network \citep{friedman1982victims}, 2) by including consultation with friends/family we tap into network focused research but are still far from it, 3) we can guess about structure of network but we don't know anything about content -- if anything, this makes my estimation conservative. 


% The RCVS even features a question that allows research in this direction --- the main reason I do not venture this way is that the order of events is unknown, and can not be assumed as securely as in the case of talking to friends and reporting to the police. This is a data limitation, which means a survey limitation, which means it can be overcome by changing the instrument. By this I do not mean the RCVS specifically --- on the contrary, in my view, any victimization survey would benefit from a question about help-seeking options taken and the order of waypoints. The technical realization of it seems to be trivial. 



***



\cite{ruback1984social}:
\textit{Research by Tyler ( 1980, Study 2) and by Skogan and Maxfield (1981) shows that learning about the victimization of friends and neighbors is associated with higher levels of fear and feelings of greater vulnerability.}

***

\textit{TBA: \citep{pescosolido1998social}
	Briefly stated, the availability of social networks (eg, size, range) determines the potential influence of social networks to affect use. But their \textbf{content} (eg, support, beliefs, values) aligns the direction of the "push"toward or away from formal medical care providers. Previous studies, which often have tapped into samples of different social classes where the nature of the network content differs, simply mapped how much influence the community had without regard to the nature of the advice likely to be given. It is not surprising studies have reported opposite  findings.}

\textit{However, the recognition of the wide variety of informal advisors has been incorporated into models of service utilization as social networks in the Health Belief Model,32 as part of the social structure in the Socio-Behavioral Model,33 and as normative beliefs in the Theory of Reasoned Action.'}

\textit{The NEM (1) reconceptualizes the dependent variable and (2) reorganizes the correlates of use, most importantly specifying the interaction of the structure of social networks (eg, number of contacts) with their content (eg, the nature of advice  and support offered). First, rather than seeing help-seeking as a single, yes-or-no, one-time decision, it conceptualizes use as patterns and sequences of these contacts. Second, when illness occurs, it is dealt with in the context of an a priori set of community supporters who provide (or fail to provide) assistance with problems. This set can be large or small, knowledgeable or ignorant of medical options, helpful or unsupportive. Depending on what kind of social networks individuals have, they will be predisposed to talk to informal advisors, seek professional care, or see a need for other social services. Network size calibrates the degree of the "push/pull"into or out of care; it does not, by itself, dictate formal use. What these networks hold in terms of experience, attitudes, and knowledge create the trajectory. For example, a large social network means a reserve of people to assist in informal care and to offer a good deal of social pressure.}

\textit{In the same way, having social networks that an individuals can rely on does not mean that such networks will be consulted for any particular problem. In the network tradition, this difference is referenced by the terms "latent"social networks, which refers to the possible contacts an individual has, and "activated" social networks, which refers to social network members who are actually contacted on any occasion. TODO: Look into people who has dense communication with relatives but did not inform them about victimization.}

!!! \textit{NEM argues that simply adding network variables to a fairly long list of costs and benefits under a traditional model mixes different theoretical understandings of the process of seeking care.8}

\textit{The finding from the usual analysis, that severity is the prime influence, may be comforting in its appeal to common sense and in the reassurance that those with the worst problems do receive care. However, with community-based care emphasizing early intervention and cost containment, this finding is of little comfort to providers or administrators. There is limited policy and treatment value in having contact with persons only when they are very ill, particularly because we do not know who will continue in care or why.}

***

***

\cite{thoits2011mechanisms}:
\textit{Although family and friends intend these acts to be helpful, they can be ineffective for two broad reasons. First, primary group members are typically upset about the individual’s situation themselves, even more so if it creates serious disruptions in their own lives. They are invested in the problem being resolved as quickly as possible to alleviate their own and their loved one’s distress. Invested supporters therefore may minimize the threatening aspects of the problem, insist on maintaining a positive outlook, or pressure the person to recover or problem solve before he or she is ready.}

\textit{In short, because primary group members are often emotionally invested in the person’s recovery and because they frequently are unfamiliar with the specific demands of the stressor, their information, advice, appraisals, and encouragement are likely to be relatively ineffective at softening its health and mental health consequences. More effective coping assistance should come instead from experientially similar others in the secondary network, to whom we now turn.}

\textit{How does the distressed person come to know that his or her secondary ties include experientially similar others? Social networks spread information about their members’ lives, especially information about tragedies and triumphs. The individual may already know of a similar other, a similar other may hear the bad news and voluntarily come forward, or people in the network may put the two in touch. This is Granovetter’s (1973) point about the “strength of weak ties”: the more extensive the person’s primary and secondary group networksand the more extensive their networks—the higher the probability the distressed individual will encounter someone with a past similar experience (Gottlieb 1985).}

\textit{Due to past experience, similar others can provide coping assistance that is closely tailored to the exigencies of the problematic situation and known to be effective in solving problems and dampening upset (Suitor and Pillemer 2000; Suitor, Pillemer, and Keeton 1995; Thoits 1986b), lessening situational demands and negative affect directly. Furthermore, guidance and encouragement from knowledgeable others should enhance the distressed individual’s sense of control and reinforce his or her active coping efforts (Berkman 1995). The advice, feedback, guidance, and encouragement supplied by experienced others should be more efficacious in buffering the impacts of adversity than the well-intentioned but less wellinformed coping assistance offered by members of the person’s primary network.}

\textit{By empirically tracing the various pathways from  social relationships or supportive acts to physical and emotional well-being, researchers will more closely target potential points of intervention. To date, most interventions have been atheoretical, ad hoc in design, and hit or miss in their effects, producing mixed and disappointing findings (Cohen 2004; Cohen et al. 2000; Hogan, Linden, and Najarian 2002). A clear understanding of how social relationships work to protect well-being is necessary (Gottlieb 2000).}

***


\cite{smith2008social}: 
\textit{health-related—ends.  Traditionally, most studies of social network effects on health actually focused on a related phenomenon, social support, a conflation that continues to a large extent today despite calls for change (Berkman \& Glass 2000). Owing predominantly to data and methodological limitations, early studies operationalized social networks as an individual-level measure of the number of social contacts a person has (structural support, or its quantitative aspect) or how helpful they are, as subjectively reported by the person (functional support, or its qualitative aspect).}

\textit{Rather than focusing, as social support studies do, on the (mere) existence of ties and conceptualizing network features as a trait reducible to the level of the individual (that is, person A has X level of social support, whereas person B has Y level of social support—a form of egocentric reduction of network information), social network studies actually map subjects’ networks and probe the impact of particular network components and kinds of ties.2 Rather than simply count the number of people in an individual’s social network (e.g., person A has N friends and person B has M friends) or rate their relative helpfulness, social network studies examine their interrelationships by focusing explicitly on the specific network links. As such, social network studies involve}

***

I do not observe victim-advisor and offender-advisor relationship \citep{knoth2019reporting}, but if anything this should make my estimation more conservative. 


%\textit{Add on Matsueda, Wikstrom, and Coleman's boat) \citep{mcgloin2011crime, matsueda2017toward}? Guess, no.
	
	
	Organizational context \citep{goudriaan2004reporting, garofalo1987school, collins1987job}
	
	\cite{knoth2019reporting}
	\textit{Crimes are least likely to be reported if the offender is part of the same triad as the victim and the advisor, and crimes are most likely to be reported if the victim, the advisor, and the offender are all strangers. Study 1 also found that males are more likely to protect friends who are offenders (by advising against reporting), while females are more likely to protect friends who are victims (by advising them to report). Study 2 found that the effect of these relationships on reporting is conditioned by the nature of the organization to which the offender belongs, such that males are particularly likely to protect their friends in athletic organizations and fraternities when accused of minor property crimes. Both studies found that gender differences in the advice to report are moderated by characteristics of the crime and triad structure.
	}





% \footnote{This is not the only regard in which fraud stands out. While in other crime types younger victims (18--24) are less likely to report than the middle aged (35--49), in case of fraud it is the other way round: young victims are the most likely reporters --- with the elders (65+) being the least likely to notify the police. The latter, by itself, should not come as a surprise: it is known that the elderly may be unwilling to report fraud \citep{parti2023if}. However, one of the expected reasons to conceal fraud victimization is feeling of guilt, and fear to compromise one's autonomy in the eyes of others -- but in my modeling there is no statistically significant association of self-blame with reporting, and neither with informing friends or family (Table \ref{table:regr_told}) -- which again singles fraud out among other crime types. I am inclined to believe that these discrepancies may be due to varying nature of fraud in different age groups, but this issue requires further investigation.}.
	
	
	
	
	\cite{pescosolido1992beyond}:\textit{ illness triggers a dynamic, social process of coping. It is through contact with others that individuals deal with situations of medical uncertainty and find ways to solve emotional and physical problems (Kulka, Veroff, and Douvan 1979). Individuals in social networks are more than an influence on help seeking, they are caregivers and advisors, part of a "therapy managing group" (Janzen 1978).}
	
	\cite{goudriaan2004reporting}: \textit{According to Gartner and Macmillan (1995), however, the unpredicted finding might be in part due to measurement error in the ICVS. Contact crimes by intimates or acquaintances are less likely to be mentioned in victimization surveys than similar crimes by unknown offenders because such surveys often do not encourage respondents to think about intimate violence as relevant for the survey (Cantor \& Lynch, 2000; Kalish, 1974). Victims of intimate violence tend to mention only the more serious crimes that they have also reported to the police. Put another way, victims of less serious attacks appear more willing to reveal these in surveys if the offender was a stranger. This tendency will bias estimates of the effects of victimoffender relationship on reporting to the police (Gartner \& Macmillan, 1995).}
	
	%		\section*{Conclusion}
	%		\label{sec:conclusion}